Les travaux du comité mettent en évidence une réalité centrale : \textbf{les \AssociationsRegionalesIHEDN{} disposent probablement d'un potentiel significatif de contribution à la résilience nationale, mais celui-ci demeure aujourd'hui insuffisamment identifié}.

La spécificité des auditeurs réside dans la diversité de leurs parcours, leur culture stratégique commune et leur capacité à produire des analyses transversales. Cette valeur ajoutée ne relève pas d'une substitution aux autorités publiques, ni d'une gestion directe des crises par l'association.

Dans ce contexte, les \AssociationsRegionalesIHEDN{} peuvent agir principalement avant la crise, par la préparation, la sensibilisation et la mise en relation des compétences, puis après la crise, par l'analyse, la production de bilans de retour d'expérience et la diffusion des enseignements. Pendant l'événement, leur rôle collectif demeure nécessairement limité par les cadres juridiques, hiérarchiques et opérationnels existants.

La mise en place d'un questionnaire constitue, à cet égard, une étape déterminante. L'analyse de l'annuaire 2022 montre en effet que le nombre d'auditeurs recensés, le nombre d'adhérents et l'engagement effectif ne se recouvrent pas. Elle confirme que les associations ne peuvent fonder leur positionnement sur une simple estimation de vivier : elles doivent disposer de données actualisées, qualifiées et interprétées avec prudence.

Le questionnaire permettrait ainsi de passer d'une intuition générale sur la richesse du réseau à une connaissance objectivée des profils, des disponibilités, des domaines d'expertise et de la motivation à contribuer. Son intérêt n'est pas de présumer une mobilisation opérationnelle, mais de clarifier ce que les \AssociationsRegionalesIHEDN{} peuvent réellement apporter, dans quelles temporalités et sous quelles conditions institutionnelles.

Plus largement, ce travail ouvre une réflexion sur la place des citoyens engagés face aux crises contemporaines. Dans un contexte marqué par l'incertitude, une hybridation des menaces et des vulnérabilités informationnelles, la résilience nationale repose de plus en plus sur une orchestration institutionnelle, territoriale et sociétale.

Les \AssociationsRegionalesIHEDN{}, par leur nature et par leur réseau, peuvent y contribuer de manière crédible, en valorisant les compétences de leurs auditeurs, en renforçant la culture de résilience et en mettant leur capacité d'analyse, de préparation et de capitalisation au service de l'intérêt général.